\documentclass[11pt]{article}
\usepackage[utf8]{inputenc}
\usepackage[lmargin=2cm, tmargin=3cm, rmargin=2cm, bmargin=2cm, headheight=5cm]{geometry}
\usepackage{graphicx}
\usepackage{subfigure}
\usepackage{enumerate}
\usepackage{listings}
\usepackage{caption}
\usepackage{float}
\usepackage{hyperref}
\usepackage{fancyhdr}
\usepackage{amsmath}

\pagestyle{fancy}{
\rhead{\includegraphics[width=30mm]{../other/logo_ist.png}}
\lhead{TCFE lab report}
}


\begin{document}
\thispagestyle{empty}
\begin{figure}[h]
	\centering
	\includegraphics[width=.6\textwidth]{../../figlib/IST_A_CMYK_POS.pdf}
\end{figure}

\begin{center}
	\huge{Circuit Theory and Electronic Fundamentals}
	
	\huge{Lab Report - T1}
	
	\vspace{30pt}
	
	\large{Professor: José Sousa}
	
	\vspace{20pt}
	
	\large{95801 - João Domingos}
	
	\large{96382 - Francisco Cadavez}
	
	\large{97087 - Miguel Fernandes}
\end{center}

\pagebreak

\tableofcontents

\pagebreak

\section{Introduction}
\hspace{12pt} In this laboratorial session we analyzed an electrical circuit composed of four meshes, eight nodes and eleven total branches: 7 resistors (R1-R7), 1 voltage source (Va), 1 current source (Id), 1 voltage controlled current source (Ib) and, finally, 1 current controlled voltage source (Vc).

\begin{figure}[h]
	\centering
	\includegraphics[width=.5\textwidth]{Circuit_diagram.pdf}
	\caption{Circuit diagram}
	\label{fig:circuit}
\end{figure}

In section 2 we applied both the mesh and nodal methods, together with mathematical support from an octave script to solve the resulting systems of linear equations, analized the circuit theoretically. In section \ref{simulation}, after recreating the circuit using an Ngspice script and using an operating point analysis we arrived at the values for the currents in each mesh and the voltages in each node. 
Finally, in section \ref{conclusion}, we used both of the previous results (theoretical and simulated) and by comparing them reach the conclusion that there is a perfect agreement.

Below is the representation of the directions and names of the currents and nodes used throughout the assignment, along with the location of the GND that has been used in the simulation and the theoretical analysis.
\begin{figure}[h]
	\centering
	\includegraphics[width=.4\textwidth]{Circuit_specs.pdf}
	\caption{Node names and current directions}
	\label{fig:circuit_spec}
\end{figure}

\pagebreak


\pagebreak
\section{Theoretical Analysis}

\subsection{Mesh Analysis Method}
\hspace{12pt} The circuit is composed of 4 meshes. The mesh analysis method tells us that the current inside a mesh is constant, and therefore, we can use Kirchhoff's Voltage Law in each mesh to determine the current.

The four elemental meshes in the circuits are: mesh A, composed by the nodes 0, B, C and H; mesh B, composed by the nodes C, D, E and H; mesh C, composed by the nodes 0, H, G and F; and mesh D, composed by the nodes H, E and F. We define $I_A$, $I_B$, $I_C$ and $I_D$ as the currents inside meshes A, B, C and D, respectively. 

We will now analyze each individual mesh, using Kirchhoff's Voltage Law.
\vspace{10pt}

Mesh A:
\begin{equation}
    \begin{split}
        R_1I_A+R_3(I_A+I_B)+R_4(I_A+I_C)=V_a \iff \\
        \iff (R_1+R_3+R_4)I_A+R_3I_B+R_4I_C=V_a
    \end{split}
\end{equation}

Mesh B:
\begin{equation}
    \begin{split}
        I_B=K_bV_B \iff I_B=K_bR_3(I_A+I_B) \iff \\
        \iff K_bR_3I_A+(K_BR_3-1)I_B=0
    \end{split}
\end{equation}

Mesh C:
\begin{equation}
    \begin{split}
        R_4(I_A+I_C)+R_6I_C+R_7I_C+V_C=0 \iff \\
        \iff R_4I_A+(R_4+R_6+R_7+K_c)I_C=0
    \end{split}
\end{equation}

Mesh D:
\begin{equation}
    I_D=I_d  \hspace{30pt} \textit{(by observation)}
\end{equation}

\vspace{20pt}
With the 4 equations above, we can form the following linear equation system:
\vspace{20pt}

$\begin{bmatrix}
    R_1 + R_3 + R_4 & R_3        & R_4             & 0 \\
    K_bR_3          & K_bR_3 - 1 & 0               & 0 \\
    R_4             & 0          & R_4+R_6+R_7-K_c & 0 \\
    0               & 0          & 0               & 1 
\end{bmatrix}$ $\cdot$
$\begin{bmatrix}
     I_A \\ I_B \\I_C \\I_D
\end{bmatrix}$ =
$\begin{bmatrix}
    V_a \\ 0 \\ 0 \\ I_d
\end{bmatrix}$

\vspace{20pt}
Using Octave to calculate the solution of the system we obtain:
\vspace{20pt}

\begin{figure}[h]
	\begin{center}
	    \begin{minipage}{.3\textwidth}
		\flushright
		\begin{tabular}{|c|}
		    \hline
		    $I_A$ \\
		    \hline
		    $I_B$ \\
		    \hline
		    $I_C$ \\
		    \hline
		    $I_D$ \\
		    \hline
		\end{tabular}
	    \end{minipage}
	    \hspace{-8pt}
	    \begin{minipage}{.3\textwidth}
		\flushleft
		\input{currents.tex}
	    \end{minipage}
	\end{center}
	\caption{Current values for the different meshes}
	\label{theory_currents}
\end{figure}

 Knowing the currents in all the circuit meshes, we can use Ohm's Law in each resistor to calculate the voltage drop and thus extrapolate the values for all of the node voltages.
 \newpage
 


\subsection{Node Analysis Method}
\hspace{12pt} This circuit is composed of 8 nodes, which have been named in the image in the beginning of the document (Figure ~\ref{fig:circuit_spec}). The Kirchhoff's Current Law says that the sum of the current that enters a node, with the current that leaves a node, must be zero. However this law cannot be applied in nodes connected to voltage sources, so we can only use this law in 4 of the 8 nodes.

We will now analyze the 4 nodes using Kirchhoff's Current Law.


\vspace{10pt}
Node C:
\begin{gather}
    I_1+I_2-I_3=0 \iff \nonumber \\ 
    \iff G_1(V_B-V_C)+G_2(V_D-V_C)-G_3(V_C-V_H)=0 \iff \\
    \iff G_1V_B-(G_1+G_2+G_3)V_C+G_2V_D+G_3V_H \nonumber
\end{gather}

Node D:
\begin{gather}
    -I_2+I_b=0 \iff \nonumber \\
    \iff -G_2(V_D-V_C)+K_b(V_C-V_H) \iff \\
    \iff (K_b+G_2)V_C-G_2V_D-K_bV_H=0 \nonumber
\end{gather}

Node E:
\begin{gather}
    -I_b-I_5+I_d=0 \iff \nonumber \\
    \iff -K_b(V_C-V_H)-G_5(V_E-V_H)+I_d=0 \iff \\
    \iff K_bV_C+G_5V_E-(K_b+G_5)V_H=I_d \nonumber
\end{gather}

Node G:
\begin{gather}
    I_c-I_7=0 \iff \nonumber \\
    \iff G_6(0-V_G)-G_7(V_G-V_F)=0 \iff \\
    \iff -(G_6+G_7)V_G+G_7V_F=0 \nonumber
\end{gather}
\vspace{5pt}

Assuming node A has 0 potential, we can conclude by observation:
\begin{equation}
    V_B=V_a
\end{equation}
\vspace{5pt}

The voltage controlled voltage source also gives:
\begin{gather}
    V_c=K_cI_c \iff \nonumber \\
    \iff V_H-V_F=K_cG_6(0-V_G) \iff \\
    \iff -V_F+K_cG_6V_G+V_H=0 \nonumber
\end{gather}
\vspace{5pt}

We get the last equation using the Ohm's Law in resistor 4:
\begin{gather}
    V_H-0=R_4(I_A+I_C) \iff \nonumber \\
    \iff V_H=R_4(G1(V_B-V_C)+G_6(0-V_G)) \iff \\
    \iff R_4G_1V_B-R_4G_1V_C-R_4G_6V_G-V_H=0 \nonumber
\end{gather}

Using the 7 equations above, we have the following linear equations system:
\vspace{10pt}

\fontsize{8}{12}\selectfont
$\begin{bmatrix}
    G1 & -(G1+G2+G3) & G2 & 0 & 0 & 0 & G3 \\
    0 & Kb+G2 & -G2 & 0 & 0 & 0 & -Kb \\
    0 & Kb & 0 & G5 & 0 & 0 & -(Kb+G5) \\
    0 & 0 & 0 & 0 & G7 & -(G6+G7) & 0 \\
    1 & 0 & 0 & 0 & 0 & 0 & 0 \\
    0 & 0 & 0 & 0 & -1 & KcG6 & 1 \\
    R4G1 & -R4G1 & 0 & 0 & 0 & -R4G6 & -1
\end{bmatrix}$ $\cdot$
$\begin{bmatrix}
     V_B \\V_C \\V_D \\V_E \\V_F \\V_G \\V_H
\end{bmatrix}$ =
$\begin{bmatrix}
    0 \\0 \\I_d \\0 \\V_a \\0 \\0
\end{bmatrix}$

\fontsize{11}{12}\selectfont
\vspace{20pt}
Using Octave to calculate the solution of the system we obtain:

\begin{figure}[h]
	\begin{center}
	    \begin{minipage}{.3\textwidth}
		\flushright
		\begin{tabular}{|c|}
		    \hline
		    $V_A$ \\
		    \hline
		    $V_B$ \\
		    \hline
		    $V_C$ \\
		    \hline
		    $V_D$ \\
		    \hline
		    $V_E$ \\
		    \hline
		    $V_F$ \\
		    \hline
		    $V_G$ \\
		    \hline
		    $V_H$ \\
		    \hline
		\end{tabular}
	    \end{minipage}
	    \hspace{-8pt}
	    \begin{minipage}{.3\textwidth}
		\flushleft
		\input{voltages.tex}
	    \end{minipage}
	\end{center}
	\caption{Electric potential of the 8 nodes}
	\label{theory_voltages}
\end{figure}

We can now use Ohm's Law to calculate the current flowing through each resistor which, consequently, allows us to reach the values for the currents in the 4 meshes.


\subsection{Method Comparison}
\hspace{12pt} We have now calculated both the currents and the voltages of the entire circuit however, having used 2 different methods does not guarantee us they agree with each other. In order to test said agreement we will analyze 2 resistors: $R_1$ and $R_5$. \vspace{12pt}

Using Ohm's Law on $R_1$ (and following the current convention we set) we hope to see that $V_B - V_C = R_1I_A$ and by carrying out said calculation we get that $V_B - V_C = 1.93585V$ and $R_1I_A = 1.93585V$. \vspace{12pt}

Let us now use the same method to analyze R5, since we are assuming the current in $R_5$ flows torward node $V_H$, we should see that $V_E - V_H = R_5(I_D - I_B)$. $V_E - V_H = 9.00903V$, $R_5(I_D - I_B) = 9.00907V$ this difference most likely comes from rounding errors within the Octave program but it correspondes to an immesurable error so we can say with certainty that both the nodal and mesh analysis methods agree with each other's results.

\pagebreak
\section{Simulation}
\label{simulation}
Using the software Ngspice, we recreated the circuit (Figure ~\ref{fig:circuit}) with the specifications as seen in Figure ~\ref{fig:circuit_spec} and after performing an \textbf{operational point analysis} (.op in the ngspice script) obtained the folowing results:

\begin{figure}[h]
	\begin{center}
		\begin{tabular}{|c|c|}
			\hline
			\textbf{Name} & \textbf{Value (A or V)} \\
			\hline
			\input{op_tab.tex}
		\end{tabular}
	\end{center}
	\caption{Current and voltage values obtained from the ngspice simulation}
	\label{sim_results}
\end{figure}




\pagebreak
\section{Conclusion}
\label{Conclusion}
\hspace{12pt} Now that we've obtained all the results from both the theoretical and the simulation analysis, we will compare them ...
\subsection{$V_S$ constant and non zero}
\hspace{12pt} Performing nodal analysis (octave table on the left) and operating point analysis (ngspice table on the right) for a dc voltage source $V_S$, we obtain the following results (Figure \ref{comp1}) in which we can clearly see the only difference being matters of rounding and the number of significant figures the Octave table provides.

\begin{figure}[h]
	\begin{minipage}{.3\textwidth}
		\flushright
		\begin{tabular}{|c|}
		    \hline
		    $V_1$ \\
		    \hline
		    $V_2$ \\
		    \hline
		    $V_3$ \\
		    \hline
		    $V_4$ \\
		    \hline
		    $V_5$ \\
		    \hline
		    $V_6$ \\
		    \hline
		    $V_7$ \\
		    \hline
		    $V_8$ \\
		    \hline
		\end{tabular}
	\end{minipage}	
	\hspace{-8pt}
	\begin{minipage}{.3\textwidth}
		\flushleft
		\input{voltages1.tex}
	\end{minipage}
	\hspace{10pt}
	\begin{minipage}{.3\textwidth}
		\begin{tabular}{|c|c|}
			\hline
			\textbf{Name} & \textbf{Value (A or V)} \\
			\hline
			\input{op_1_tab.tex}
		\end{tabular}
	\end{minipage}
	\caption{Theoretical results using Octave (left) and simulation results from Ngspice (right)}    
	\label{comp1}
\end{figure}
\pagebreak
\subsection{$V_S = 0$ and capacitor replaced by voltage source}
\hspace{12pt} After setting the voltage source $V_S$ off and replacing the capacitor by a dc voltage source $V_x = V_6 - V_8$ as obtained by the methods above (this is used in order to simulate the charged capacitor) and performing the same analysis as in Figure \ref{comp1}, we arrive at results in Figure \ref{comp2}. In the Ngspice table, we can see a repeated pattern of voltages whose order of magnite is around $10^{-15}$, this can be explained as approximations that the software makes while solving the respective systems os linear equations. Nonetheless, these voltages are so incredibly small that they are negligible. 

\begin{figure}[h]
	\begin{minipage}{.3\textwidth}
		\flushright
		\begin{tabular}{|c|}
		    \hline
		    $V_1$ \\
		    \hline
		    $V_2$ \\
		    \hline
		    $V_3$ \\
		    \hline
		    $V_4$ \\
		    \hline
		    $V_5$ \\
		    \hline
		    $V_6$ \\
		    \hline
		    $V_7$ \\
		    \hline
		    $V_8$ \\
		    \hline
		\end{tabular}
	\end{minipage}	
	\hspace{-8pt}
	\begin{minipage}{.3\textwidth}
		\flushleft
		\input{voltages2.tex}
	\end{minipage}
	\hspace{10pt}
	\begin{minipage}{.3\textwidth}
		\begin{tabular}{|c|c|}
			\hline
			\textbf{Name} & \textbf{Value (A or V)} \\
			\hline
			\input{op_2_tab.tex}
		\end{tabular}
	\end{minipage}
	\caption{Theoretical results using Octave (left) and simulation results from Ngspice (right)}    
	\label{comp2}
\end{figure}

We can thus see that in both the analysed cases, the voltage results from both methods agree with each other within a margin of $10^{-14} V$, we can then conclude that the Ngspice software does appear to use the same methods as the ones introduced in the theoretical analysis section.


\end{document}
